\section{Scope, Evidence, and Limits}

\subsection{Question, range, and terminal date}

This account asks how a small Swiss work of priestly formation became a large, durable body of postconciliar Catholic traditionalism outside regular canonical structures; why successive Roman negotiations failed; how each side understood the conflict; and what changed when the Society consecrated bishops without pontifical mandate for a second time on 1 July 2026. It begins with Lefebvre's earlier missionary and conciliar career and ends with the Dicastery for the Doctrine of the Faith's decree and explanatory note of 2 July 2026, together with the Society's initial response. Sources and mutable status were checked through 16 July 2026.

The geography is international but has a Swiss center. Fribourg supplied the first house and diocesan erection; Écône in the Diocese of Sion became the principal seminary and the site of both sets of episcopal consecrations; Menzingen became the general house. France, Germany, the United States, Argentina, Australia, and missions elsewhere show how a controversy rooted in Western Europe's postconciliar transition became a worldwide organization.

The reader-facing series name, ``Traditional Priestly Institutes,'' is only a bibliographic grouping for connected histories. It is not a canonical classification. This account uses ``Society'' and ``SSPX'' as conventional English names for the \emph{Fraternitas Sacerdotalis Sancti Pii X}; it does not imply that the body presently possesses the status of a society of apostolic life in the Catholic Church.

\subsection{Evidence layers and viewpoint control}

\begin{threestable}{Layer}{Representative evidence}{Use and limit}
Contemporary institutional record & Erection and suppression documents, correspondence, protocols, decrees, statutes, election notices, and dated communiqués & Establishes an act, proposal, or participant position. A party's document is not a neutral judgment on the whole dispute.\\
Later institutional memory & SSPX, Holy See, FSSP, diocesan, and journalistic retrospectives & Establishes how institutions remember and frame their past. It must be checked against contemporary acts and may contain advocacy or compressed chronology.\\
Juridical and magisterial act & Papal letters, canonical decrees, the Code of Canon Law, and dicastery notes & Establishes the authority's operative judgment or rule at the stated date. It does not settle every historical motive or reproduce the Society's counterargument.\\
Modern reconstruction & Historical, sociological, and canonical scholarship & Supplies context, comparison, and interpretation. Its claims are separated from official status and from participant testimony.\\
\end{threestable}

The publication therefore uses paired attribution at points of dispute. ``The Holy See declared'' reports an official act. ``The Society argued'' reports the SSPX's own defense. Neither formula is a typographic verdict. Where the Holy See's 2026 act is the present official source, the account states that result plainly and then records the Society's rejection as a response, not as a cancellation of the act.

\subsection{A bounded canonical module}

Historical reconstruction governs the work. Current-status passages additionally follow the universal canon-law source rules: the governing body is the Latin Church's 1983 Code of Canon Law as revised, especially Book VI effective 8 December 2021; the promulgating authority is the Roman Pontiff; special papal and dicastery acts are named; and the cutoff is 16 July 2026. The work does not advise anyone whether to approach a minister, incur a penalty, contest an act, or seek sacramental remedies. Those are pastoral and canonical questions for competent authority and the facts of an individual case.

\evidenceframe{Dicastery for the Doctrine of the Faith, decree and explanatory note, 2 July 2026.}{The decree declared specified bishops excommunicated after the 1 July consecrations; the explanatory note declared SSPX sacred ministers schismatic and subject to the penalty in canon 1364 §1, preserved the 1996 case-by-case test for lay formal adherence, and stated consequences for sacramental ministry from that point forward.}{The Italian acts are reported by paraphrase, not translated by the project. They do not authorize this study to determine an individual's culpability or remedy.}

This discipline is especially important because popular descriptions---``approved,'' ``suppressed,'' ``excommunicated,'' ``in full communion,'' ``valid,'' or ``traditional Mass''---often collapse distinct persons, acts, and periods. The history assigns each term a date, an authority, and a defined object.

Latin, French, and Italian records are paraphrased; the project supplies no unofficial translation of them. The publication reproduces no extended third-party prose, official translation, archival image, or facsimile. Short attributed phrases and source-identifying titles---including ``Catholic Rome,'' ``neo-Modernist and neo-Protestant tendencies,'' ``eternal Rome,'' ``Operation Survival,'' and ``auxiliary bishops without jurisdiction''---are retained only where needed to analyze participant positions. The work otherwise uses dates, names, bibliographic facts, and source-grounded paraphrase. The research audit records the full claim-to-source alignment, negative findings, access record, and rights decisions.

\subsection{Publication boundary}

This source-audited historical study was checked through 16 July 2026. It reports ecclesiastical acts and participant claims at their own dates; it is not a canonical opinion, a pastoral directive, or an adjudication of contested theological arguments. Independent historical, theological, and canonical review remains outstanding.
